\section{从账本看早明的地方接收}

本节按账本的时间序列观察地区、财政、边疆与社会问题。官修实录、诏令和奏疏强于保存中央选择记录的行动与日期；地方志、碑刻、契约、工程遗存及其他政权材料才可能校正具体地点的执行、损失与反应。因此，以下叙事不把行政分类、报送数字或动机判断直接等同于地方社会本身。

\subsection{户籍、田亩与军户}

\eventanchor{ML-1368-001}{明·1368（洪武元年）一月18日}
\paragraph{1368（洪武元年）一月18日：红巾起义：胡美攻占福州} 账本首先限定我们能够确认的范围：编年候选的年序可由官修与后出材料交叉；具体月日、参与者、战果或数字必须回查所列材料，不能将单一叙事当作定论。可供互读的材料包括《太祖实录》洪武元年条；《明史·兵志》及相关列传；边镇、地方志或对方材料。这些文本并非同一份现场证言，而是从朝廷、地方或跨政权的位置留下观察。事件之所以进入行政链，与；它带来的可追踪变化是它改变了后续调兵、补给或地方秩序的条件；战果和伤亡须分开核对。译写候选以编年材料定位；实录记载反映朝廷选择，崇祯期须另以奏疏、地方及敌对政权材料交叉。

\eventanchor{ML-1368-002}{明·1368（洪武元年）一月23日}
\paragraph{1368（洪武元年）一月23日：朱元璋自立为明朝洪武皇帝（注意明清使用年号而不是庙名）} 账本首先限定我们能够确认的范围：编年候选的年序可由官修与后出材料交叉；具体月日、参与者、战果或数字必须回查所列材料，不能将单一叙事当作定论。可供互读的材料包括《太祖实录》洪武元年条；《明史·本纪》及相关列传；诏令、奏疏或会典版本。这些文本并非同一份现场证言，而是从朝廷、地方或跨政权的位置留下观察。事件之所以进入行政链，与；它带来的可追踪变化是它影响后续人事与政策议程；动机和派系标签不能由单一叙事裁定。译写候选以编年材料定位；实录记载反映朝廷选择，崇祯期须另以奏疏、地方及敌对政权材料交叉。

\eventanchor{ML-1368-003}{明·1368（洪武元年）二月17日}
\paragraph{1368（洪武元年）二月17日：明军攻克福建，俘获陈有定，并处决} 账本首先限定我们能够确认的范围：编年候选的年序可由官修与后出材料交叉；具体月日、参与者、战果或数字必须回查所列材料，不能将单一叙事当作定论。可供互读的材料包括《太祖实录》洪武元年条；《明史·兵志》及相关列传；边镇、地方志或对方材料。这些文本并非同一份现场证言，而是从朝廷、地方或跨政权的位置留下观察。事件之所以进入行政链，与；它带来的可追踪变化是它改变了后续调兵、补给或地方秩序的条件；战果和伤亡须分开核对。译写候选以编年材料定位；实录记载反映朝廷选择，崇祯期须另以奏疏、地方及敌对政权材料交叉。

\eventanchor{ML-1368-004}{明·1368（洪武元年）三月1日}
\paragraph{1368（洪武元年）三月1日：明军攻克山东} 账本首先限定我们能够确认的范围：编年候选的年序可由官修与后出材料交叉；具体月日、参与者、战果或数字必须回查所列材料，不能将单一叙事当作定论。可供互读的材料包括《太祖实录》洪武元年条；《明史·兵志》及相关列传；边镇、地方志或对方材料。这些文本并非同一份现场证言，而是从朝廷、地方或跨政权的位置留下观察。事件之所以进入行政链，与；它带来的可追踪变化是它改变了后续调兵、补给或地方秩序的条件；战果和伤亡须分开核对。译写候选以编年材料定位；实录记载反映朝廷选择，崇祯期须另以奏疏、地方及敌对政权材料交叉。

\eventanchor{ML-1368-005}{明·1368（洪武元年）四月16日}
\paragraph{1368（洪武元年）四月16日：明军攻克开封} 账本首先限定我们能够确认的范围：编年候选的年序可由官修与后出材料交叉；具体月日、参与者、战果或数字必须回查所列材料，不能将单一叙事当作定论。可供互读的材料包括《太祖实录》洪武元年条；《明史·兵志》及相关列传；边镇、地方志或对方材料。这些文本并非同一份现场证言，而是从朝廷、地方或跨政权的位置留下观察。事件之所以进入行政链，与；它带来的可追踪变化是它改变了后续调兵、补给或地方秩序的条件；战果和伤亡须分开核对。译写候选以编年材料定位；实录记载反映朝廷选择，崇祯期须另以奏疏、地方及敌对政权材料交叉。

\eventanchor{ML-1368-006}{明·1368（洪武元年）四月18日}
\paragraph{1368（洪武元年）四月18日：明军到达广州并接受何震的投降} 账本首先限定我们能够确认的范围：编年候选的年序可由官修与后出材料交叉；具体月日、参与者、战果或数字必须回查所列材料，不能将单一叙事当作定论。可供互读的材料包括《太祖实录》洪武元年条；《明史·兵志》及相关列传；边镇、地方志或对方材料。这些文本并非同一份现场证言，而是从朝廷、地方或跨政权的位置留下观察。事件之所以进入行政链，与；它带来的可追踪变化是它改变了后续调兵、补给或地方秩序的条件；战果和伤亡须分开核对。译写候选以编年材料定位；实录记载反映朝廷选择，崇祯期须另以奏疏、地方及敌对政权材料交叉。

\subsection{地方接收与边地调发}

\eventanchor{ML-1368-007}{明·1368（洪武元年）四月25日}
\paragraph{1368（洪武元年）四月25日：明军击败阔克帖木儿并占领洛阳} 账本首先限定我们能够确认的范围：编年候选的年序可由官修与后出材料交叉；具体月日、参与者、战果或数字必须回查所列材料，不能将单一叙事当作定论。可供互读的材料包括《太祖实录》洪武元年条；《明史·兵志》及相关列传；边镇、地方志或对方材料。这些文本并非同一份现场证言，而是从朝廷、地方或跨政权的位置留下观察。事件之所以进入行政链，与；它带来的可追踪变化是它改变了后续调兵、补给或地方秩序的条件；战果和伤亡须分开核对。译写候选以编年材料定位；实录记载反映朝廷选择，崇祯期须另以奏疏、地方及敌对政权材料交叉。

\eventanchor{ML-1368-008}{明·1368（洪武元年）五月26日}
\paragraph{1368（洪武元年）五月26日：明军攻克梧州} 账本首先限定我们能够确认的范围：编年候选的年序可由官修与后出材料交叉；具体月日、参与者、战果或数字必须回查所列材料，不能将单一叙事当作定论。可供互读的材料包括《太祖实录》洪武元年条；《明史·兵志》及相关列传；边镇、地方志或对方材料。这些文本并非同一份现场证言，而是从朝廷、地方或跨政权的位置留下观察。事件之所以进入行政链，与；它带来的可追踪变化是它改变了后续调兵、补给或地方秩序的条件；战果和伤亡须分开核对。译写候选以编年材料定位；实录记载反映朝廷选择，崇祯期须另以奏疏、地方及敌对政权材料交叉。

\eventanchor{ML-1368-009}{明·1368（洪武元年）七月}
\paragraph{1368（洪武元年）七月：明军攻克广西} 账本首先限定我们能够确认的范围：编年候选的年序可由官修与后出材料交叉；具体月日、参与者、战果或数字必须回查所列材料，不能将单一叙事当作定论。可供互读的材料包括《太祖实录》洪武元年条；《明史·兵志》及相关列传；边镇、地方志或对方材料。这些文本并非同一份现场证言，而是从朝廷、地方或跨政权的位置留下观察。事件之所以进入行政链，与；它带来的可追踪变化是它改变了后续调兵、补给或地方秩序的条件；战果和伤亡须分开核对。译写候选以编年材料定位；实录记载反映朝廷选择，崇祯期须另以奏疏、地方及敌对政权材料交叉。

\eventanchor{ML-1368-010}{明·1368（洪武元年）九月20日}
\paragraph{1368（洪武元年）九月20日：明军占领大都（改称北平），元朝逃往内蒙古；元朝就这样结束了} 账本首先限定我们能够确认的范围：编年候选的年序可由官修与后出材料交叉；具体月日、参与者、战果或数字必须回查所列材料，不能将单一叙事当作定论。可供互读的材料包括《太祖实录》洪武元年条；《明史·兵志》及相关列传；边镇、地方志或对方材料。这些文本并非同一份现场证言，而是从朝廷、地方或跨政权的位置留下观察。事件之所以进入行政链，与；它带来的可追踪变化是它改变了后续调兵、补给或地方秩序的条件；战果和伤亡须分开核对。译写候选以编年材料定位；实录记载反映朝廷选择，崇祯期须另以奏疏、地方及敌对政权材料交叉。

\eventanchor{ML-1368-011}{明·1368（洪武元年）十一月}
\paragraph{1368（洪武元年）十一月：明军攻克保定} 账本首先限定我们能够确认的范围：编年候选的年序可由官修与后出材料交叉；具体月日、参与者、战果或数字必须回查所列材料，不能将单一叙事当作定论。可供互读的材料包括《太祖实录》洪武元年条；《明史·兵志》及相关列传；边镇、地方志或对方材料。这些文本并非同一份现场证言，而是从朝廷、地方或跨政权的位置留下观察。事件之所以进入行政链，与；它带来的可追踪变化是它改变了后续调兵、补给或地方秩序的条件；战果和伤亡须分开核对。译写候选以编年材料定位；实录记载反映朝廷选择，崇祯期须另以奏疏、地方及敌对政权材料交叉。

\eventanchor{ML-1368-012}{明·1368（洪武元年）十二月26日}
\paragraph{1368（洪武元年）十二月26日：明军攻克赵州} 账本首先限定我们能够确认的范围：编年候选的年序可由官修与后出材料交叉；具体月日、参与者、战果或数字必须回查所列材料，不能将单一叙事当作定论。可供互读的材料包括《太祖实录》洪武元年条；《明史·兵志》及相关列传；边镇、地方志或对方材料。这些文本并非同一份现场证言，而是从朝廷、地方或跨政权的位置留下观察。事件之所以进入行政链，与；它带来的可追踪变化是它改变了后续调兵、补给或地方秩序的条件；战果和伤亡须分开核对。译写候选以编年材料定位；实录记载反映朝廷选择，崇祯期须另以奏疏、地方及敌对政权材料交叉。

\subsection{海路、使节与港口}

\eventanchor{ML-1368-013}{明·1368（洪武元年）十二月}
\paragraph{1368（洪武元年）十二月：明军攻克平定} 账本首先限定我们能够确认的范围：编年候选的年序可由官修与后出材料交叉；具体月日、参与者、战果或数字必须回查所列材料，不能将单一叙事当作定论。可供互读的材料包括《太祖实录》洪武元年条；《明史·兵志》及相关列传；边镇、地方志或对方材料。这些文本并非同一份现场证言，而是从朝廷、地方或跨政权的位置留下观察。事件之所以进入行政链，与；它带来的可追踪变化是它改变了后续调兵、补给或地方秩序的条件；战果和伤亡须分开核对。译写候选以编年材料定位；实录记载反映朝廷选择，崇祯期须另以奏疏、地方及敌对政权材料交叉。

\eventanchor{ML-1368-014}{明·1368（洪武元年）}
\paragraph{1368（洪武元年）：卧虎炮为明军所用} 账本首先限定我们能够确认的范围：编年候选的年序可由官修与后出材料交叉；具体月日、参与者、战果或数字必须回查所列材料，不能将单一叙事当作定论。可供互读的材料包括《太祖实录》洪武元年条；《明史·选举志》或《艺文志》；文集、碑刻或书目材料。这些文本并非同一份现场证言，而是从朝廷、地方或跨政权的位置留下观察。事件之所以进入行政链，与；它带来的可追踪变化是它提供制度和传播史的锚点，不能据此推断社会整体接受。译写候选以编年材料定位；实录记载反映朝廷选择，崇祯期须另以奏疏、地方及敌对政权材料交叉。

\eventanchor{ML-1368-015}{明·1368（洪武元年）}
\paragraph{1368（洪武元年）：国子监创建} 账本首先限定我们能够确认的范围：编年候选的年序可由官修与后出材料交叉；具体月日、参与者、战果或数字必须回查所列材料，不能将单一叙事当作定论。可供互读的材料包括《太祖实录》洪武元年条；《明史·本纪》及相关列传；诏令、奏疏或会典版本。这些文本并非同一份现场证言，而是从朝廷、地方或跨政权的位置留下观察。事件之所以进入行政链，与；它带来的可追踪变化是它影响后续人事与政策议程；动机和派系标签不能由单一叙事裁定。译写候选以编年材料定位；实录记载反映朝廷选择，崇祯期须另以奏疏、地方及敌对政权材料交叉。

\eventanchor{ML-1369-001}{明·1369（洪武二年）一月9日}
\paragraph{1369（洪武二年）一月9日：明军攻克太原} 账本首先限定我们能够确认的范围：编年候选的年序可由官修与后出材料交叉；具体月日、参与者、战果或数字必须回查所列材料，不能将单一叙事当作定论。可供互读的材料包括《太祖实录》洪武二年条；《明史·兵志》及相关列传；边镇、地方志或对方材料。这些文本并非同一份现场证言，而是从朝廷、地方或跨政权的位置留下观察。事件之所以进入行政链，与；它带来的可追踪变化是它改变了后续调兵、补给或地方秩序的条件；战果和伤亡须分开核对。译写候选以编年材料定位；实录记载反映朝廷选择，崇祯期须另以奏疏、地方及敌对政权材料交叉。

\eventanchor{ML-1369-002}{明·1369（洪武二年）三月3日}
\paragraph{1369（洪武二年）三月3日：明军攻克大同} 账本首先限定我们能够确认的范围：编年候选的年序可由官修与后出材料交叉；具体月日、参与者、战果或数字必须回查所列材料，不能将单一叙事当作定论。可供互读的材料包括《太祖实录》洪武二年条；《明史·兵志》及相关列传；边镇、地方志或对方材料。这些文本并非同一份现场证言，而是从朝廷、地方或跨政权的位置留下观察。事件之所以进入行政链，与；它带来的可追踪变化是它改变了后续调兵、补给或地方秩序的条件；战果和伤亡须分开核对。译写候选以编年材料定位；实录记载反映朝廷选择，崇祯期须另以奏疏、地方及敌对政权材料交叉。

\eventanchor{ML-1369-003}{明·1369（洪武二年）三月}
\paragraph{1369（洪武二年）三月：宋濂、王邑开始编修《元史》} 账本首先限定我们能够确认的范围：编年候选的年序可由官修与后出材料交叉；具体月日、参与者、战果或数字必须回查所列材料，不能将单一叙事当作定论。可供互读的材料包括《太祖实录》洪武二年条；《明史·本纪》及相关列传；诏令、奏疏或会典版本。这些文本并非同一份现场证言，而是从朝廷、地方或跨政权的位置留下观察。事件之所以进入行政链，与；它带来的可追踪变化是它影响后续人事与政策议程；动机和派系标签不能由单一叙事裁定。译写候选以编年材料定位；实录记载反映朝廷选择，崇祯期须另以奏疏、地方及敌对政权材料交叉。

\subsection{都城、工程与转运}

\eventanchor{ML-1369-004}{明·1369（洪武二年）四月18日}
\paragraph{1369（洪武二年）四月18日：明军攻克山西，李思齐逃往临洮} 账本首先限定我们能够确认的范围：编年候选的年序可由官修与后出材料交叉；具体月日、参与者、战果或数字必须回查所列材料，不能将单一叙事当作定论。可供互读的材料包括《太祖实录》洪武二年条；《明史·兵志》及相关列传；边镇、地方志或对方材料。这些文本并非同一份现场证言，而是从朝廷、地方或跨政权的位置留下观察。事件之所以进入行政链，与；它带来的可追踪变化是它改变了后续调兵、补给或地方秩序的条件；战果和伤亡须分开核对。译写候选以编年材料定位；实录记载反映朝廷选择，崇祯期须另以奏疏、地方及敌对政权材料交叉。

\eventanchor{ML-1369-005}{明·1369（洪武二年）五月21日}
\paragraph{1369（洪武二年）五月21日：李思齐向明军投降} 账本首先限定我们能够确认的范围：编年候选的年序可由官修与后出材料交叉；具体月日、参与者、战果或数字必须回查所列材料，不能将单一叙事当作定论。可供互读的材料包括《太祖实录》洪武二年条；《明史·兵志》及相关列传；边镇、地方志或对方材料。这些文本并非同一份现场证言，而是从朝廷、地方或跨政权的位置留下观察。事件之所以进入行政链，与；它带来的可追踪变化是它改变了后续调兵、补给或地方秩序的条件；战果和伤亡须分开核对。译写候选以编年材料定位；实录记载反映朝廷选择，崇祯期须另以奏疏、地方及敌对政权材料交叉。

\eventanchor{ML-1369-006}{明·1369（洪武二年）五月23日}
\paragraph{1369（洪武二年）五月23日：明军攻克兰州} 账本首先限定我们能够确认的范围：编年候选的年序可由官修与后出材料交叉；具体月日、参与者、战果或数字必须回查所列材料，不能将单一叙事当作定论。可供互读的材料包括《太祖实录》洪武二年条；《明史·兵志》及相关列传；边镇、地方志或对方材料。这些文本并非同一份现场证言，而是从朝廷、地方或跨政权的位置留下观察。事件之所以进入行政链，与；它带来的可追踪变化是它改变了后续调兵、补给或地方秩序的条件；战果和伤亡须分开核对。译写候选以编年材料定位；实录记载反映朝廷选择，崇祯期须另以奏疏、地方及敌对政权材料交叉。

\eventanchor{ML-1369-007}{明·1369（洪武二年）六月8日}
\paragraph{1369（洪武二年）六月8日：明军攻克平凉} 账本首先限定我们能够确认的范围：编年候选的年序可由官修与后出材料交叉；具体月日、参与者、战果或数字必须回查所列材料，不能将单一叙事当作定论。可供互读的材料包括《太祖实录》洪武二年条；《明史·兵志》及相关列传；边镇、地方志或对方材料。这些文本并非同一份现场证言，而是从朝廷、地方或跨政权的位置留下观察。事件之所以进入行政链，与；它带来的可追踪变化是它改变了后续调兵、补给或地方秩序的条件；战果和伤亡须分开核对。译写候选以编年材料定位；实录记载反映朝廷选择，崇祯期须另以奏疏、地方及敌对政权材料交叉。

\eventanchor{ML-1369-008}{明·1369（洪武二年）七月20日}
\paragraph{1369（洪武二年）七月20日：明军攻克上都} 账本首先限定我们能够确认的范围：编年候选的年序可由官修与后出材料交叉；具体月日、参与者、战果或数字必须回查所列材料，不能将单一叙事当作定论。可供互读的材料包括《太祖实录》洪武二年条；《明史·兵志》及相关列传；边镇、地方志或对方材料。这些文本并非同一份现场证言，而是从朝廷、地方或跨政权的位置留下观察。事件之所以进入行政链，与；它带来的可追踪变化是它改变了后续调兵、补给或地方秩序的条件；战果和伤亡须分开核对。译写候选以编年材料定位；实录记载反映朝廷选择，崇祯期须另以奏疏、地方及敌对政权材料交叉。

\eventanchor{ML-1369-009}{明·1369（洪武二年）九月22日}
\paragraph{1369（洪武二年）九月22日：明军攻克庆阳} 账本首先限定我们能够确认的范围：编年候选的年序可由官修与后出材料交叉；具体月日、参与者、战果或数字必须回查所列材料，不能将单一叙事当作定论。可供互读的材料包括《太祖实录》洪武二年条；《明史·兵志》及相关列传；边镇、地方志或对方材料。这些文本并非同一份现场证言，而是从朝廷、地方或跨政权的位置留下观察。事件之所以进入行政链，与；它带来的可追踪变化是它改变了后续调兵、补给或地方秩序的条件；战果和伤亡须分开核对。译写候选以编年材料定位；实录记载反映朝廷选择，崇祯期须另以奏疏、地方及敌对政权材料交叉。
